\section{IEEE 33 节点数值实验}\label{sec:experiments}
\subsection{IEEE 33 节点配置与求解方法}
实验使用 MATPOWER 的 case33bw~\cite{matpower33}，保留 32 条有效径向支路，基准电压 12.66 kV、基准容量 10 MVA，原始有功、无功负荷同时乘以 0.65。光伏逆变器设置于 9、18、24、33 号节点，每台标称有功注入为 0.25 MW、右半圆有功--无功扰动半径为 0.125 MVA、视在功率额定值为 1 MVA。额外新增的单位功率因数负荷全部接入 18 号节点。电压限制为 $0.95\le V_i\le1.05$ pu，线路电流统一采用 $I_l^{\max}=0.31$ pu。本算例采用统一线路电流限值 $0.31$ pu，四个逆变器增益的搜索范围为 $[0,4]^4$。

已有 Python 原型对每一负荷单元分别计算非线性状态包络、分块 Hessian 电压界、复功率几何电流界和有功--无功耦合功率圆上界，并自适应细分到从零到目标容量的全部区间均通过测试。本文采用 Python/SciPy 数值潮流求解器，使用 Jacobian、分块 Hessian 和复功率几何界检查候选负荷区间。结果属于采用所述物理模型和双精度数值实现的\emph{计算容量}；并用非线性潮流结果展示两种增益配置的电压响应。

\begin{figure}[htbp]
\centering
\begin{tikzpicture}
\begin{axis}[
width=.88\linewidth,height=6.15cm,
ybar,bar width=19pt,ymin=0.95,ymax=1.18,
symbolic x coords={统一增益,前期协调,改进协调},
xtick=data,
ylabel={候选安全接入容量（MW）},
ymajorgrids=true,grid style={gray!20},
nodes near coords={\pgfmathprintnumber[fixed,precision=4]{\pgfplotspointmeta}},
nodes near coords align={vertical},
every node near coord/.append style={font=\small},
enlarge x limits=.20,
axis line style={gray!70},
tick label style={font=\small},
ylabel style={font=\small},
]
\addplot[draw=teal!70!black,fill=teal!32] coordinates {
(统一增益,1.0586) (前期协调,1.1363) (改进协调,1.1500)};
\end{axis}
\end{tikzpicture}
\caption{IEEE 33 节点不同逆变器增益下的安全接入容量计算结果（网络配置及运行约束相同）。}
\label{fig:capacity}
\end{figure}


\begin{table}[ht]
\centering
\caption{IEEE 33 节点三组逆变器增益下的计算接入容量与负荷单元数}
\label{tab:results}
\begin{tabular}{@{}p{32mm}p{72mm}rr@{}}
\toprule
设计方案 & 增益（节点 9、18、24、33） & 容量/MW & 单元数\\
\midrule
统一增益 & $(1.92,1.92,1.92,1.92)$ & 1.0586 & 58\\
协调增益（前期） & $(3.70,1.92,2.12,2.52)$ & 1.1363 & 62\\
协调增益（改进） & $(3.80105,1.89598,2.39520,4.00000)$ & 1.1500 & 60\\
\bottomrule
\end{tabular}
\end{table}

相比统一增益，协调设计的计算接入容量增加
\[
1.1500-1.0586=0.0914\ {\rm MW}=91.4\ {\rm kW},
\qquad \frac{0.0914}{1.0586}\approx8.63\%.
\]
同一网络和设备配置下，协调增益能够提高算法得到的接入容量，并且可直接输出四台逆变器的控制参数。

\subsection{逆变器协调对电压的改善}
为观察控制参数的影响，在四台光伏保持同一测试扰动 $w_{p,j}=0$、$w_{q,j}=-0.0125$ pu 时，用非线性潮流求解器计算节点 18 在不同新增负荷下的电压。图~\ref{fig:voltage} 比较统一增益与协调增益的响应：在新增负荷为 $1.00$ MW 时，18 号节点电压分别约为 $0.95348$ pu 和 $0.95673$ pu。协调增益提高了电压水平，从而为接入新增负荷留出更大的运行裕度。

\begin{figure}[htbp]
\centering
\begin{tikzpicture}
\begin{axis}[
width=.90\linewidth,height=6.3cm,
xmin=0,xmax=1.20,ymin=.943,ymax=.993,
xlabel={18 号节点新增负荷（MW）},
ylabel={18 号节点电压（pu）},
grid=both,major grid style={gray!18},
legend style={font=\small,at={(.03,.03)},anchor=south west,fill=white},
tick label style={font=\small},
label style={font=\small},
]
\addplot[blue!80!black,mark=*,mark size=1.1pt,thick] coordinates {
(0,.98739798)(.2,.98110375)(.4,.97458404)(.6,.96782096)
(.8,.96079377)(1,.95347815)(1.05,.95160091)(1.1,.94970338)
(1.15,.94778502)(1.187,.94635170)};
\addlegendentry{统一增益 $g_j=1.92$}
\addplot[teal!80!black,mark=square*,mark size=1.15pt,thick] coordinates {
(0,.98936283)(.2,.98327888)(.4,.97699040)(.6,.97048222)
(.8,.96373687)(1,.95673412)(1.05,.95494053)(1.1,.95312898)
(1.15,.95129905)(1.187,.94993281)};
\addlegendentry{异构增益 $(3.801,1.896,2.395,4.0)$}
\addplot[red!75!black,dashed,thick] coordinates {(0,.95)(1.20,.95)};
\addlegendentry{最低电压 $0.95$ pu}
\end{axis}
\end{tikzpicture}
\caption{实际非线性物理限幅潮流求解器输出的负荷--电压曲线：四台光伏均采用一个可行扰动 $w_{p,j}=0,\ w_{q,j}=-0.0125$ pu。该图是\emph{固定扰动工况}比较，不代表已搜索到全局最坏扰动。}
\label{fig:voltage}
\end{figure}


\subsection{同一 IEEE 33 节点模型下的算法消融}

为分析结构化上界对容量计算的贡献，我们在相同馈线、相同光伏不确定集、相同逆变器增益和\emph{相同完整视在功率圆约束}下，对比两种早期标量电流上界基线（整体 Hessian、分块 Hessian）与本文采用的复功率几何电流方法。本文算法采用复功率几何电流界，两个标量 Hessian 方法作为容量计算基线。三种方法使用一致的电压、光伏不确定性和逆变器视在功率约束。

取角度分辨率 $N_\theta=2048$，初始负荷单元宽度 $0.02$ MW，最小宽度 $0.00025$ MW。表~\ref{tab:ablation} 给出协调增益方案下的结果。表中“通过”表示算法接受目标容量，“未通过”表示该方法的上界计算未满足接入条件。

\begin{table}[htbp]
\centering
\caption{协调增益 $(3.80105,1.89598,2.39520,4.0)$ 下各方法的容量区间检查}
\label{tab:ablation}
\begin{tabular}{@{}p{65mm}ccc@{}}
\toprule
上界计算方案 & $0.80$ MW & $1.00$ MW & $1.15$ MW\\
\midrule
整体 Hessian + 标量电流界 & 通过 & 未通过 & 未通过\\
分块 Hessian + 标量电流界 & 通过 & 未通过 & 未通过\\
分块 Hessian + 复功率几何电流界 & 通过 & 通过 & 通过\\
\bottomrule
\end{tabular}
\end{table}

两个历史标量基线的首次无法接受区间分别位于约 $0.8319$ MW 和 $0.8329$ MW，主导原因均为线路电流界过于保守。改用复功率几何界后，同样的增益与约束下可以覆盖整个 $[0,1.15]$ MW。这说明，\emph{充分利用电流的复功率结构比单纯改变 Hessian 范数更有价值}。对于统一增益，标量电流方法在约 $1.031$ MW 停止接受，而复功率几何方法可通过 $1.0586$ MW，趋势一致。

\subsection{运行时间与计算量}
在三种对照方法均能够通过的目标 $L=0.80$ MW 下，使用完全相同的算例与精度参数，在同一 Python 数值环境中各连续运行三次，表~\ref{tab:runtime} 给出墙钟耗时中位数。相较标量 Hessian 方法，复功率几何搜索增加了角度计算，但明显改善了高容量下的约束判定能力。正式算法不再重复求取标量电流上界。

\begin{table}[htbp]
\centering
\caption{协调增益下 $L=0.80$ MW 的耗时对照（同环境三次中位数）}
\label{tab:runtime}
\begin{tabular}{@{}p{80mm}rr@{}}
\toprule
上界计算方案 & 时间/s & 单元求值次数\\
\midrule
整体 Hessian + 标量电流界 & 0.604 & 46\\
分块 Hessian + 标量电流界 & 0.699 & 40\\
分块 Hessian + 复功率几何电流界 & 1.162 & 40\\
\bottomrule
\end{tabular}
\end{table}

对较高的 $L=1.15$ MW，复功率几何法评估 $62$ 个负荷单元，一次运行耗时约 $1.78$ s；两个历史标量基线在约 $0.83$ MW 就提前停止，因此比较运行时间时同时考虑能够通过的容量范围。所有对照方法均使用一致的光伏完整视在功率约束。

\subsection{角度分辨率敏感性}
复功率角度最大化采用有限角度网格实现，角度数 $N_\theta$ 决定离散修正项 $\pi K_l/N_\theta$ 的大小。保持协调增益及目标容量 $L=1.15$ MW 不变，表~\ref{tab:angular} 列出角度分辨率对容量检查和运行时间的影响。256 个角度时，线路电流上界使约 $1.138$ MW 的单元未能通过；取 512 个及以上角度则均接受完整目标容量。这说明角度搜索既是计算成本来源，也直接影响容量上界的紧度。

\begin{table}[htbp]
\centering
\caption{协调增益下目标 $L=1.15$ MW 的角度网格敏感性}
\label{tab:angular}
\begin{tabular}{@{}rrrr@{}}
\toprule
角度数 $N_\theta$ & 容量检查 & 单元求值数 & 单次耗时/s\\
\midrule
256 & 未通过 & 73 & 1.425\\
512 & 通过 & 66 & 1.311\\
1024 & 通过 & 62 & 1.376\\
2048 & 通过 & 62 & 1.619\\
4096 & 通过 & 62 & 2.162\\
\bottomrule
\end{tabular}
\end{table}

\section{结论}
本文针对含光伏不确定性及逆变器 Volt/VAR 调节的径向配电网，提出安全接入容量计算与协调增益联合设计算法。方法以非线性 DistFlow 求解器为基础，利用分块 Hessian 表示节点电压的非线性变化，再由半圆支撑函数和复功率角度界分别处理光伏扰动及线路电流，通过完整负荷区间检查得到安全接入容量，并据此更新逆变器控制增益。

IEEE 33 节点实验中，统一下垂增益方案的计算容量为 $1.0586$ MW，协调增益方案为 $1.1500$ MW，提高约 $8.63\%$。同模型消融表明，标量 Hessian 电流上界约在 $0.83$ MW 因保守性而提前限制容量，复功率几何电流界则允许完整通过 $1.15$ MW 的容量检查。复功率几何方法增加了一定角度搜索开销，但显著提高了容量估计的利用率；因此后续容量设计统一采用这一种线路电流上界。

因此，对不同物理量分别采用匹配其结构的灵敏度与几何上界，并与逆变器增益更新结合，是一种直接、可实现的配电网安全新增负荷容量设计方法。
